%%%
%
% $Autor: Wings $
% $Date: 2024-10-31 11:15:45Z $
% $File: TensorFlowLiteGoogleColab.tex $
% $Version: 1 $
%
%%%


\chapter{TensorFlow Lite Object Detection API in Google Colab}
	
	An object detection model is trained to detect the presence and location of multiple classes of objects. For example, a model might be trained with images that contain various pieces of fruit, along with a label that specifies the class of fruit they represent (e.g. an apple, a banana, or a strawberry), and data specifying where each object appears in the image.
	
	When an image is subsequently provided to the model, it will output a list of the objects it detects, the location of a bounding box that contains each object, and a score that indicates the confidence that detection was correct. \cite{TensorFlowLiteObjectDetection:2024}
	
	\section{Introduction}
	
	In this practice, we will implement the TensorFlow Object Detection API to train an SSD-MobileNet model or EfficientDet model with a custom dataset and convert it to TensorFlow Lite format. By working through this Google Colab, we will be able to create and download a TFLite model that you can run on your PC, an Android phone or an edge device. ~\ref{CustomCoinModel} \cite{TensorFlowLiteObjectDetection:2024}
	
	\begin{figure}
		\begin{center}
			\includegraphics[width=0.7\linewidth]{TensorFlowLite/CustomCoinModel.png}
			\caption{Custom SSD-MobileNet-FPNLite model}
			\label{CustomCoinModel}
		\end{center}
	\end{figure}
	
	\subsection{Working in Collab}
	Colab provides a virtual machine in your browser complete with a Linux OS, filesystem, Python environment, and best of all, a free GPU. It comes with most TensorFlow backend requirements (like CUDA and cuDNN) pre-installed. Simply click the play button on sections of code in this notebook to execute them on the virtual machine.
	
	\section{Procedure}
	
	\subsection{Gather and Label Training Images}
	Before we start training, we need to gather and label images that will be used for training the object detection model. A good starting point for a proof-of-concept model is 200 images. The training images should have random objects in the image along with the desired objects, and should have a variety of backgrounds and lighting conditions. \cite{TensorFlowLiteObjectDetection:2024}
	
	When you've finished gathering and labeling images, you should have a folder full of images and corresponding .xml data annotation file for each image. An example of a labeled image and the image folder for my coin detector model are shown in the figure ~\ref{ImageFolder} ~\ref{LabeledImage}
	
	\begin{figure}
		\begin{center}
			\includegraphics[width=0.7\linewidth]{TensorFlowLite/LabeledImage.png}
			\caption{Labeled Image}
			\label{LabeledImage}
		\end{center}
	\end{figure}
	
	\begin{figure}
		\begin{center}
			\includegraphics[width=0.7\linewidth]{TensorFlowLite/ImageFolder.png}
			\caption{Image Folder}
			\label{ImageFolder}
		\end{center}
	\end{figure}
	
	\subsection{Install TensorFlow Object Detection Dependencies}
	
	First, we will install the TensorFlow Object Detection API in this Google Colab instance. This requires cloning the \SHELL{TensorFlow models repository} \url{https://github.com/tensorflow/models} and running a couple installation commands. We need to click the play button to run the following sections of code in the Google collab. \cite{TensorFlowLiteObjectDetection:2024}
	
	
	\begin{lstlisting}[language=bash, numbers=left, numberstyle=\tiny, stepnumber=1, numbersep=5pt, frame=single]
		# Clone the tensorflow models repository from GitHub
		!pip uninstall Cython -y # Temporary fix for "No module named 'object_detection'" error
		!git clone --depth 1 https://github.com/tensorflow/models
	\end{lstlisting}
	
	\begin{lstlisting}[language=bash, numbers=left, numberstyle=\tiny, stepnumber=1, numbersep=5pt, frame=single]
		# Copy setup files into models/research folder
		%%bash
		cd models/research/
		protoc object_detection/protos/*.proto --python_out=.
		#cp object_detection/packages/tf2/setup.py .
	\end{lstlisting}
	
	\begin{lstlisting}[language=bash, numbers=left, numberstyle=\tiny, stepnumber=1, numbersep=5pt, frame=single]
		# Modify setup.py file to install the tf-models-official repository targeted at TF v2.8.0
		import re
		with open('/content/models/research/object_detection/packages/tf2/setup.py') as f:
		s = f.read()
		
		with open('/content/models/research/setup.py', 'w') as f:
		# Set fine_tune_checkpoint path
		s = re.sub('tf-models-official>=2.5.1',
		'tf-models-official==2.8.0', s)
		f.write(s)
	\end{lstlisting}
	
	\begin{lstlisting}[language=bash, numbers=left, numberstyle=\tiny, stepnumber=1, numbersep=5pt, frame=single, caption=Python script for Installing Object Detection Dependencies]
		# Install the Object Detection API (NOTE: This block takes about 10 minutes to finish executing)
		
		# Need to do a temporary fix with PyYAML because Colab isn't able to install PyYAML v5.4.1
		!pip install pyyaml==5.3
		!pip install /content/models/research/
		
		# Need to downgrade to TF v2.8.0 due to Colab compatibility bug with TF v2.10 (as of 10/03/22)
		!pip install tensorflow==2.8.0
		
		# Install CUDA version 11.0 (to maintain compatibility with TF v2.8.0)
		!pip install tensorflow_io==0.23.1
		!wget https://developer.download.nvidia.com/compute/cuda/repos/ubuntu1804/x86_64/cuda-ubuntu1804.pin
		!mv cuda-ubuntu1804.pin /etc/apt/preferences.d/cuda-repository-pin-600
		!wget http://developer.download.nvidia.com/compute/cuda/11.0.2/local_installers/cuda-repo-ubuntu1804-11-0-local_11.0.2-450.51.05-1_amd64.deb
		!dpkg -i cuda-repo-ubuntu1804-11-0-local_11.0.2-450.51.05-1_amd64.deb
		!apt-key add /var/cuda-repo-ubuntu1804-11-0-local/7fa2af80.pub
		!apt-get update && sudo apt-get install cuda-toolkit-11-0
		!export LD_LIBRARY_PATH=/usr/local/cuda-11.0/lib64:$LD_LIBRARY_PATH
	\end{lstlisting}
	
	Now test our installation by running \SHELL{model\_builder\_tf2\_test.py} to make sure everything is working as expected. Run the following code block and confirm that it finishes without errors.
	
	\begin{lstlisting}[language=bash, numbers=left, numberstyle=\tiny, stepnumber=1, numbersep=5pt, frame=single]
		# Run Model Bulider Test file, just to verify everything's working properly
		!python /content/models/research/object_detection/builders/model_builder_tf2_test.py
	\end{lstlisting}
	
	
	\subsection{Upload Image Dataset and Prepare Training Data}
	
	In this section, we will upload our data and prepare it for training with TensorFlow. Then upload our images, split them into train, validation and test folders and then run scripts for creating TFRecords from our data.
	
	First, on your local PC, zip all your training images and XML files into a single folder called \FILE{images.zip}.
	
	There are three options for moving the image files to this Colab instance.
	
	\subsubsection{Option 1. Upload through Google Colab}
	
	Upload the \FILE{images.zip} file to the Google Colab instance by clicking the "Files" icon on the left hand side of the browser, and then the "Upload to session storage" icon. Select the zip folder to upload it. \cite{TensorFlowLiteObjectDetection:2024}
	
	\begin{figure}
		\begin{center}
			\includegraphics[width=0.7\linewidth]{TensorFlowLite/FilesCollab.png}
			\caption{Files}
			\label{FilesCollab}
		\end{center}
	\end{figure}
	
	\subsubsection{Option 2. Copy from Google Drive}
	
	You can also upload your images to your personal Google Drive, mount the drive on this Colab session, and copy them over to the Colab filesystem. This option works well if you want to upload the images beforehand so you don't have to wait for them to upload each time you restart this Colab. If you have more than 50MB worth of images, I recommend using this option.
	
	First, upload the \FILE{images.zip} file to your Google Drive, and make note of the folder you uploaded them to. Replace MyDrive/path/to/images.zip with the path to your zip file. Then, run the following block of code to mount your Google Drive to this Colab session and copy the folder to this filesystem.
	
	\subsubsection{Option 3. Use coin detection dataset}
	
	Use already existing coin dataset containing 750 labeled images of pennies, nickels, dimes, and quarters which is available in the storage cloud.
	
	\begin{lstlisting}[language=bash, numbers=left, numberstyle=\tiny, stepnumber=1, numbersep=5pt, frame=single]
		!wget -O /content/images.zip https://www.dropbox.com/s/gk57ec3v8dfuwcp/CoinPics_11NOV22.zip?dl=0  # United States coin images
	\end{lstlisting}
	
	\subsubsection{Split images into train, validation, and test folders}
	
	Now that the dataset is uploaded, unzip it and create some folders to hold the images. These directories are created in the /content folder in this instance's filesystem using the script ~\ref{unzipscript}. You can browse the filesystem by clicking the "Files" icon on the left.
	
	\begin{lstlisting}[language=bash, numbers=left, numberstyle=\tiny, stepnumber=1, numbersep=5pt, frame=single, label=unzipscript, caption={unzip dataset using python script}]
		!mkdir /content/images
		!unzip -q images.zip -d /content/images/all
		!mkdir /content/images/train; mkdir /content/images/validation; mkdir /content/images/test
	\end{lstlisting}
	
	Next, we need to split the images into train, validation, and test sets. Here's what each set is used for:
	\begin{enumerate}
		\item \SHELL{Train}: These are the actual images used to train the model. In each step of training, a batch of images from the "train" set is passed into the neural network. The network predicts classes and locations of objects in the images. The training algorithm calculates the loss (i.e. how "wrong" the predictions were) and adjusts the network weights through backpropagation.
		\item \SHELL{Validation}: Images from the "validation" set can be used by the training algorithm to check the progress of training and adjust hyperparameters (like learning rate). Unlike "train" images, these images are only used periodically during training (i.e. once every certain number of training steps).
		\item \SHELL{Test}: These images are never seen by the neural network during training. They are intended to be used by a human to perform final testing of the model to check how accurate the model is.
	\end{enumerate}
	Python script ~\ref{datasetsplit} to randomly move 80 percent of the images to the "train" folder, 10 percent to the "validation" folder, and 10 percent to the "test" folder.
	
	\begin{lstlisting}[language=bash, numbers=left, numberstyle=\tiny, stepnumber=1, numbersep=5pt, frame=single, label=datasetsplit, caption={Splitting dataset using python script}]
		!wget https://raw.githubusercontent.com/EdjeElectronics/TensorFlow-Lite-Object-Detection-on-Android-and-Raspberry-Pi/master/util_scripts/train_val_test_split.py
		!python train_val_test_split.py
	\end{lstlisting}
	
	\subsubsection{Create Labelmap and TFRecords}
	Finally, we need to create a labelmap for the detector and convert the images into a data file format called TFRecords, which are used by TensorFlow for training. We'll use Python scripts to automatically convert the data into TFRecord format. Before running them, we need to define a labelmap for our classes.
	The code section ~\ref{classes} will create a file \FILE{labelmap.txt} that contains a list of classes.
	
	\begin{lstlisting}[language=bash, numbers=left, numberstyle=\tiny, stepnumber=1, numbersep=5pt, frame=single, label=classes, caption={LabelMap}]
		### This creates a a "labelmap.txt" file with a list of classes the object detection model will detect.
		%%bash
		cat <<EOF >> /content/labelmap.txt
		penny
		nickel
		dime
		quarter
		EOF
	\end{lstlisting}
	
	The script ~\ref{TFRecord} will create TFRecord files for the train and validation datasets, as well as a file \FILE{labelmap.pbtxt} which contains the labelmap in a different format.
	
	\begin{lstlisting}[language=bash, numbers=left, numberstyle=\tiny, stepnumber=1, numbersep=5pt, frame=single, label=TFRecord, caption={TF Record files}]
		# Download data conversion scripts
		! wget https://raw.githubusercontent.com/EdjeElectronics/TensorFlow-Lite-Object-Detection-on-Android-and-Raspberry-Pi/master/util_scripts/create_csv.py
		! wget https://raw.githubusercontent.com/EdjeElectronics/TensorFlow-Lite-Object-Detection-on-Android-and-Raspberry-Pi/master/util_scripts/create_tfrecord.py
		# Create CSV data files and TFRecord files
		!python3 create_csv.py
		!python3 create_tfrecord.py --csv_input=images/train_labels.csv --labelmap=labelmap.txt --image_dir=images/train --output_path=train.tfrecord
		!python3 create_tfrecord.py --csv_input=images/validation_labels.csv --labelmap=labelmap.txt --image_dir=images/validation --output_path=val.tfrecord
	\end{lstlisting}
	
	\subsection{Setting up Training Configuration}
	
	In this section, we will set up an SSD-MobileNet model training configuration. We will specifiy which pretrained TensorFlow model we want to use from the \SHELL{TensorFlow 2 Object Detection Model Zoo} \url{https://github.com/tensorflow/models/blob/master/research/object_detection/g3doc/tf1_detection_zoo.md}. Each model also comes with a configuration file that points to file locations, sets training parameters (such as learning rate and total number of training steps), and more. We will modify the configuration file for our custom training job.
	
	The first section of code lists out some models availabe in the TF2 Model Zoo and defines some filenames that will be used later to download the model and config file. This makes it easy to manage which model you're using and to add other models to the list later.
	
	Set the "chosen\_model" variable to match the name of the model you'd like to train with. It's currently set to use the model \SHELL{ssd-mobilenet-v1-quantized}. Run the next three sections to specify and download the pretrained model file and configuration file. \cite{TensorFlowLiteObjectDetection:2024}
	
	\begin{lstlisting}[language=bash, numbers=left, numberstyle=\tiny, stepnumber=1, numbersep=5pt, frame=single, label=model, caption={TF2 Object detection zoo model}]
		# Change the chosen_model variable to deploy different models available in the TF2 object detection zoo
		chosen_model = 'ssd-mobilenet-v2-fpnlite-320'
		
		MODELS_CONFIG = {
			'ssd-mobilenet-v2': {
				'model_name': 'ssd_mobilenet_v2_320x320_coco17_tpu-8',
				'base_pipeline_file': 'ssd_mobilenet_v2_320x320_coco17_tpu-8.config',
				'pretrained_checkpoint': 'ssd_mobilenet_v2_320x320_coco17_tpu-8.tar.gz',
			},
			'efficientdet-d0': {
				'model_name': 'efficientdet_d0_coco17_tpu-32',
				'base_pipeline_file': 'ssd_efficientdet_d0_512x512_coco17_tpu-8.config',
				'pretrained_checkpoint': 'efficientdet_d0_coco17_tpu-32.tar.gz',
			},
			'ssd-mobilenet-v2-fpnlite-320': {
				'model_name': 'ssd_mobilenet_v2_fpnlite_320x320_coco17_tpu-8',
				'base_pipeline_file': 'ssd_mobilenet_v2_fpnlite_320x320_coco17_tpu-8.config',
				'pretrained_checkpoint': 'ssd_mobilenet_v2_fpnlite_320x320_coco17_tpu-8.tar.gz',
			},
			# The centernet model isn't working as of 9/10/22
			#'centernet-mobilenet-v2': {
				#    'model_name': 'centernet_mobilenetv2fpn_512x512_coco17_od',
				#    'base_pipeline_file': 'pipeline.config',
				#    'pretrained_checkpoint': 'centernet_mobilenetv2fpn_512x512_coco17_od.tar.gz',
				#}
		}
		
		model_name = MODELS_CONFIG[chosen_model]['model_name']
		pretrained_checkpoint = MODELS_CONFIG[chosen_model]['pretrained_checkpoint']
		base_pipeline_file = MODELS_CONFIG[chosen_model]['base_pipeline_file']
	\end{lstlisting}
	
	Download the pretrained model file and configuration file by running the script on the following section.
	
	\begin{lstlisting}[language=bash, numbers=left, numberstyle=\tiny, stepnumber=1, numbersep=5pt, frame=single]
		# Create "mymodel" folder for holding pre-trained weights and configuration files
		%mkdir /content/models/mymodel/
		%cd /content/models/mymodel/
		
		# Download pre-trained model weights
		import tarfile
		download_tar = 'http://download.tensorflow.org/models/object_detection/tf2/20200711/' + pretrained_checkpoint
		!wget {download_tar}
		tar = tarfile.open(pretrained_checkpoint)
		tar.extractall()
		tar.close()
		
		# Download training configuration file for model
		download_config = 'https://raw.githubusercontent.com/tensorflow/models/master/research/object_detection/configs/tf2/' + base_pipeline_file
		!wget {download_config}
	\end{lstlisting}
	
	Now that we have downloaded our model and config file, we need to modify the configuration file with some high-level training parameters. The following variables are used to control training steps:
	
	\begin{enumerate}
		\item \SHELL{num\_steps}: The total amount of steps to use for training the model. A good number to start with is 40,000 steps. You can use more steps if you notice the loss metrics are still decreasing by the time training finishes. The more steps, the longer the training. Training can also be stopped early if loss flattens out before reaching the specified number of steps.
		\item \SHELL{batch\_size}: The number of images to use per training step. A larger batch size allows a model to be trained in fewer steps, but the size is limited by the GPU memory available for training. With the GPUs used in Colab instances, 16 is a good number for SSD models and 4 is good for EfficientDet models.
	\end{enumerate}
	
	Other training information, like the location of the pretrained model file, the config file, and total number of classes are also assigned in this step.
	
	
	\begin{lstlisting}[language=bash, numbers=left, numberstyle=\tiny, stepnumber=1, numbersep=5pt, frame=single]
		# Set training parameters for the model
		num_steps = 40000
		if chosen_model == 'efficientdet-d0':
		batch_size = 4
		else:
		batch_size = 16
	\end{lstlisting}
	
	\begin{lstlisting}[language=bash, numbers=left, numberstyle=\tiny, stepnumber=1, numbersep=5pt, frame=single]
		# Set file locations and get number of classes for config file
		pipeline_fname = '/content/models/mymodel/' + base_pipeline_file
		fine_tune_checkpoint = '/content/models/mymodel/' + model_name + '/checkpoint/ckpt-0'
		
		def get_num_classes(pbtxt_fname):
		from object_detection.utils import label_map_util
		label_map = label_map_util.load_labelmap(pbtxt_fname)
		categories = label_map_util.convert_label_map_to_categories(
		label_map, max_num_classes=90, use_display_name=True)
		category_index = label_map_util.create_category_index(categories)
		return len(category_index.keys())
		num_classes = get_num_classes(label_map_pbtxt_fname)
		print('Total classes:', num_classes)
	\end{lstlisting}
	
	Finally, we will rewrite the config file to use the training parameters we just specified. The following section of code will automatically replace the necessary parameters in the downloaded .config file and save it as our custom file \FILE{pipeline\_file.config}.
	
	\begin{lstlisting}[language=bash, numbers=left, numberstyle=\tiny, stepnumber=1, numbersep=5pt, frame=single]
		# Create custom configuration file by writing the dataset, model checkpoint, and training parameters into the base pipeline file
		import re
		
		%cd /content/models/mymodel
		print('writing custom configuration file')
		
		with open(pipeline_fname) as f:
		s = f.read()
		with open('pipeline_file.config', 'w') as f:
		
		# Set fine_tune_checkpoint path
		s = re.sub('fine_tune_checkpoint: ".*?"',
		'fine_tune_checkpoint: "{}"'.format(fine_tune_checkpoint), s)
		
		# Set tfrecord files for train and test datasets
		s = re.sub(
		'(input_path: ".*?)(PATH_TO_BE_CONFIGURED/train)(.*?")', 'input_path: "{}"'.format(train_record_fname), s)
		s = re.sub(
		'(input_path: ".*?)(PATH_TO_BE_CONFIGURED/val)(.*?")', 'input_path: "{}"'.format(val_record_fname), s)
		
		# Set label_map_path
		s = re.sub(
		'label_map_path: ".*?"', 'label_map_path: "{}"'.format(label_map_pbtxt_fname), s)
		
		# Set batch_size
		s = re.sub('batch_size: [0-9]+',
		'batch_size: {}'.format(batch_size), s)
		
		# Set training steps, num_steps
		s = re.sub('num_steps: [0-9]+',
		'num_steps: {}'.format(num_steps), s)
		
		# Set number of classes num_classes
		s = re.sub('num_classes: [0-9]+',
		'num_classes: {}'.format(num_classes), s)
		
		# Change fine-tune checkpoint type from "classification" to "detection"
		s = re.sub(
		'fine_tune_checkpoint_type: "classification"', 'fine_tune_checkpoint_type: "{}"'.format('detection'), s)
		
		# If using ssd-mobilenet-v2, reduce learning rate (because it's too high in the default config file)
		if chosen_model == 'ssd-mobilenet-v2':
		s = re.sub('learning_rate_base: .8',
		'learning_rate_base: .08', s)
		
		s = re.sub('warmup_learning_rate: 0.13333',
		'warmup_learning_rate: .026666', s)
		
		# If using efficientdet-d0, use fixed_shape_resizer instead of keep_aspect_ratio_resizer (because it isn't supported by TFLite)
		if chosen_model == 'efficientdet-d0':
		s = re.sub('keep_aspect_ratio_resizer', 'fixed_shape_resizer', s)
		s = re.sub('pad_to_max_dimension: true', '', s)
		s = re.sub('min_dimension', 'height', s)
		s = re.sub('max_dimension', 'width', s)
		
		f.write(s)
		
	\end{lstlisting}
	
	\subsection{Train Custom TFLite Detection Model}
	
	Before we start training, we need to load a TensorBoard session to monitor training progress. Run the following section of code, and a TensorBoard session will appear in the browser. It won't show anything yet, because we haven't started training. Once training starts, come back and click the refresh button to see the model's overall loss.
	
	\begin{lstlisting}[language=bash, numbers=left, numberstyle=\tiny, stepnumber=1, numbersep=5pt, frame=single]
		%load_ext tensorboard
		%tensorboard --logdir '/content/training/train'
	\end{lstlisting}
	
	Model training is performed using the \FILE{model\_main\_tf2.py} script from the TF Object Detection API. Training will take anywhere from 2 to 6 hours, depending on the model, batch size, and number of training steps. We've already defined all the parameters and arguments used by \FILE{model\_main\_tf2.py} in previous sections of this Colab. 
	
	\begin{lstlisting}[language=bash, numbers=left, numberstyle=\tiny, stepnumber=1, numbersep=5pt, frame=single, caption={Run Model using custom dataset}]
		# Run training!
		!python /content/models/research/object_detection/model_main_tf2.py \
		--pipeline_config_path={pipeline_file} \
		--model_dir={model_dir} \
		--alsologtostderr \
		--num_train_steps={num_steps} \
		--sample_1_of_n_eval_examples=1
	\end{lstlisting}
	
	\subsection{Convert Model to TensorFlow Lite}
	
	The model is all trained up and ready to be used for detecting objects. First, we need to export the model graph (a file that contains information about the architecture and weights) to a TensorFlow Lite-compatible format. We'll do this using the script \FILE{export\_tflite\_graph\_tf2.py}.
	
	\begin{lstlisting}[language=bash, numbers=left, numberstyle=\tiny, stepnumber=1, numbersep=5pt, frame=single]
		# Make a directory to store the trained TFLite model
		!mkdir /content/custom_model_lite
		output_directory = '/content/custom_model_lite'
		
		# Path to training directory (the conversion script automatically chooses the highest checkpoint file)
		last_model_path = '/content/training'
		
		!python /content/models/research/object_detection/export_tflite_graph_tf2.py \
		--trained_checkpoint_dir {last_model_path} \
		--output_directory {output_directory} \
		--pipeline_config_path {pipeline_file}
		
	\end{lstlisting}
	
	Next, we will take the exported graph and use the TFLiteConverter module to convert it to FlatBuffer format \FILE{.tflite}.
	
	\begin{lstlisting}[language=bash, numbers=left, numberstyle=\tiny, stepnumber=1, numbersep=5pt, frame=single]
		# Convert exported graph file into TFLite model file
		import tensorflow as tf
		
		converter = tf.lite.TFLiteConverter.from_saved_model('/content/custom_model_lite/saved_model')
		tflite_model = converter.convert()
		
		with open('/content/custom_model_lite/detect.tflite', 'wb') as f:
		f.write(tflite_model)
	\end{lstlisting}
	
	\subsection{Test TensorFlow Lite Model and Calculate mAP}
	
	We have trained our custom model and converted it to TFLite format. But how well does it actually perform at detecting objects in images? This is where the images we set aside in the test folder come in. The model never saw any test images during training, so its performance on these images should be representative of how it will perform on new images from the field.
	
	\subsubsection{Inference test images}
	
	The following code defines a function to run inference on test images. It loads the images, loads the model and labelmap, runs the model on each image, and displays the result. It also optionally saves detection results as text files so we can use them to calculate model mAP score.
	This code is based off the script \FILE{TFLite\_detection\_image.py}.
	
	
	\begin{lstlisting}[language=bash, numbers=left, numberstyle=\tiny, stepnumber=1, numbersep=5pt, frame=single]
		# Script to run custom TFLite model on test images to detect objects
		# Source: https://github.com/EdjeElectronics/TensorFlow-Lite-Object-Detection-on-Android-and-Raspberry-Pi/blob/master/TFLite_detection_image.py
		
		# Import packages
		import os
		import cv2
		import numpy as np
		import sys
		import glob
		import random
		import importlib.util
		from tensorflow.lite.python.interpreter import Interpreter
		
		import matplotlib
		import matplotlib.pyplot as plt
		
		%matplotlib inline
		
		### Define function for inferencing with TFLite model and displaying results
		
		def tflite_detect_images(modelpath, imgpath, lblpath, min_conf=0.5, num_test_images=10, savepath='/content/results', txt_only=False):
		
		# Grab filenames of all images in test folder
		images = glob.glob(imgpath + '/*.jpg') + glob.glob(imgpath + '/*.JPG') + glob.glob(imgpath + '/*.png') + glob.glob(imgpath + '/*.bmp')
		
		# Load the label map into memory
		with open(lblpath, 'r') as f:
		labels = [line.strip() for line in f.readlines()]
		
		# Load the Tensorflow Lite model into memory
		interpreter = Interpreter(model_path=modelpath)
		interpreter.allocate_tensors()
		
		# Get model details
		input_details = interpreter.get_input_details()
		output_details = interpreter.get_output_details()
		height = input_details[0]['shape'][1]
		width = input_details[0]['shape'][2]
		
		float_input = (input_details[0]['dtype'] == np.float32)
		
		input_mean = 127.5
		input_std = 127.5
		
		# Randomly select test images
		images_to_test = random.sample(images, num_test_images)
		
		# Loop over every image and perform detection
		for image_path in images_to_test:
		
		# Load image and resize to expected shape [1xHxWx3]
		image = cv2.imread(image_path)
		image_rgb = cv2.cvtColor(image, cv2.COLOR_BGR2RGB)
		imH, imW, _ = image.shape
		image_resized = cv2.resize(image_rgb, (width, height))
		input_data = np.expand_dims(image_resized, axis=0)
		
		# Normalize pixel values if using a floating model (i.e. if model is non-quantized)
		if float_input:
		input_data = (np.float32(input_data) - input_mean) / input_std
		
		# Perform the actual detection by running the model with the image as input
		interpreter.set_tensor(input_details[0]['index'],input_data)
		interpreter.invoke()
		
		# Retrieve detection results
		boxes = interpreter.get_tensor(output_details[1]['index'])[0] # Bounding box coordinates of detected objects
		classes = interpreter.get_tensor(output_details[3]['index'])[0] # Class index of detected objects
		scores = interpreter.get_tensor(output_details[0]['index'])[0] # Confidence of detected objects
		
		detections = []
		
		# Loop over all detections and draw detection box if confidence is above minimum threshold
		for i in range(len(scores)):
		if ((scores[i] > min_conf) and (scores[i] <= 1.0)):
		
		# Get bounding box coordinates and draw box
		# Interpreter can return coordinates that are outside of image dimensions, need to force them to be within image using max() and min()
		ymin = int(max(1,(boxes[i][0] * imH)))
		xmin = int(max(1,(boxes[i][1] * imW)))
		ymax = int(min(imH,(boxes[i][2] * imH)))
		xmax = int(min(imW,(boxes[i][3] * imW)))
		
		cv2.rectangle(image, (xmin,ymin), (xmax,ymax), (10, 255, 0), 2)
		
		# Draw label
		object_name = labels[int(classes[i])] # Look up object name from "labels" array using class index
		label = '%s: %d%%' % (object_name, int(scores[i]*100)) # Example: 'person: 72%'
		labelSize, baseLine = cv2.getTextSize(label, cv2.FONT_HERSHEY_SIMPLEX, 0.7, 2) # Get font size
		label_ymin = max(ymin, labelSize[1] + 10) # Make sure not to draw label too close to top of window
		cv2.rectangle(image, (xmin, label_ymin-labelSize[1]-10), (xmin+labelSize[0], label_ymin+baseLine-10), (255, 255, 255), cv2.FILLED) # Draw white box to put label text in
		cv2.putText(image, label, (xmin, label_ymin-7), cv2.FONT_HERSHEY_SIMPLEX, 0.7, (0, 0, 0), 2) # Draw label text
		
		detections.append([object_name, scores[i], xmin, ymin, xmax, ymax])
		
		
		# All the results have been drawn on the image, now display the image
		if txt_only == False: # "text_only" controls whether we want to display the image results or just save them in .txt files
		image = cv2.cvtColor(image,cv2.COLOR_BGR2RGB)
		plt.figure(figsize=(12,16))
		plt.imshow(image)
		plt.show()
		
		# Save detection results in .txt files (for calculating mAP)
		elif txt_only == True:
		
		# Get filenames and paths
		image_fn = os.path.basename(image_path)
		base_fn, ext = os.path.splitext(image_fn)
		txt_result_fn = base_fn +'.txt'
		txt_savepath = os.path.join(savepath, txt_result_fn)
		
		# Write results to text file
		# (Using format defined by https://github.com/Cartucho/mAP, which will make it easy to calculate mAP)
		with open(txt_savepath,'w') as f:
		for detection in detections:
		f.write('%s %.4f %d %d %d %d\n' % (detection[0], detection[1], detection[2], detection[3], detection[4], detection[5]))
		
		return
	\end{lstlisting}
	
	\subsubsection{Calculate mAP}
	
	Now we have a visual sense of how our model performs on test images, but how can we quantitatively measure its accuracy?
	
	One popular method for measuring object detection model accuracy is "mean average precision" (mAP). Basically, the higher the mAP score, the better your model is at detecting objects in images.
	
	\begin{lstlisting}[language=bash, numbers=left, numberstyle=\tiny, stepnumber=1, numbersep=5pt, frame=single]
		%%bash
		git clone https://github.com/Cartucho/mAP /content/mAP
		cd /content/mAP
		rm input/detection-results/*
		rm input/ground-truth/*
		rm input/images-optional/*
		wget https://raw.githubusercontent.com/EdjeElectronics/TensorFlow-Lite-Object-Detection-on-Android-and-Raspberry-Pi/master/util_scripts/calculate_map_cartucho.py
	\end{lstlisting}
	
	Next, copy the images and annotation data from the test folder to the appropriate folders inside the cloned repository. These will be used as the "ground truth data" that our model's detection results will be compared to.
	
	\begin{lstlisting}[language=bash, numbers=left, numberstyle=\tiny, stepnumber=1, numbersep=5pt, frame=single]
		!cp /content/images/test/* /content/mAP/input/images-optional # Copy images and xml files
		!mv /content/mAP/input/images-optional/*.xml /content/mAP/input/ground-truth/  # Move xml files to the appropriate folder
	\end{lstlisting}
	
	The calculator tool expects annotation data in a format that's different from the Pascal \FILE{VOC.xml} file format we're using. Fortunately, it provides an easy script, \FILE{convert\_gt\_xml.py}, for converting to the expected .txt format.
	
	\begin{lstlisting}[language=bash, numbers=left, numberstyle=\tiny, stepnumber=1, numbersep=5pt, frame=single]
		!python /content/mAP/scripts/extra/convert_gt_xml.py
	\end{lstlisting}
	
	Now we have set up the ground truth data, but we need actual detection results from our model. The detection results will be compared to the ground truth data to calculate the model's accuracy in mAP.
	
	The inference function can be used to generate detection data for all the images in the test folder. We'll use it the same as before, except this time we'll tell it to save detection results into the detection-results folder.
	
	\begin{lstlisting}[language=bash, numbers=left, numberstyle=\tiny, stepnumber=1, numbersep=5pt, frame=single]
		# Set up variables for running inference, this time to get detection results saved as .txt files
		PATH_TO_IMAGES='/content/images/test'   # Path to test images folder
		PATH_TO_MODEL='/content/custom_model_lite/detect.tflite'   # Path to .tflite model file
		PATH_TO_LABELS='/content/labelmap.txt'   # Path to labelmap.txt file
		PATH_TO_RESULTS='/content/mAP/input/detection-results' # Folder to save detection results in
		min_conf_threshold=0.1   # Confidence threshold
		
		# Use all the images in the test folder
		image_list = glob.glob(PATH_TO_IMAGES + '/*.jpg') + glob.glob(PATH_TO_IMAGES + '/*.JPG') + glob.glob(PATH_TO_IMAGES + '/*.png') + glob.glob(PATH_TO_IMAGES + '/*.bmp')
		images_to_test = min(500, len(image_list)) # If there are more than 500 images in the folder, just use 500
		
		# Tell function to just save results and not display images
		txt_only = True
		
		# Run inferencing function!
		print('Starting inference on %d images...' % images_to_test)
		tflite_detect_images(PATH_TO_MODEL, PATH_TO_IMAGES, PATH_TO_LABELS, min_conf_threshold, images_to_test, PATH_TO_RESULTS, txt_only)
		print('Finished inferencing!')
	\end{lstlisting}
	
	Finally, let's calculate mAP! One popular style for reporting mAP is the COCO metric for mAP @ 0.50:0.95. Basically, this means that mAP is calculated at several IoU thresholds between 0.50 and 0.95, and then the result from each threshold is averaged to get a final mAP score.
	The script to run the calculator tool at each IoU threshold, average the results, and report the final accuracy score. It reports mAP for each class and overall mAP. 
	
	\begin{lstlisting}[language=bash, numbers=left, numberstyle=\tiny, stepnumber=1, numbersep=5pt, frame=single]
		%cd /content/mAP
		!python calculate_map_cartucho.py --labels=/content/labelmap.txt
	\end{lstlisting}
	
	The score reported at the end is your model's overall mAP score. Ideally, it should be above 50 percent (0.50). If it isn't, you can increase your model's accuracy by adding more images to your dataset.
	
	\subsection{Deploy TensorFlow Lite Model}
	
	Now that your custom model has been trained and converted to TFLite format, it's ready to be downloaded and deployed in an application! This section shows how to download the model and provides links to instructions for deploying it on PC, or other edge devices.
	
	\subsubsection{Download TFLite model}
	
	\begin{lstlisting}[language=bash, numbers=left, numberstyle=\tiny, stepnumber=1, numbersep=5pt, frame=single]
		# Move labelmap and pipeline config files into TFLite model folder and zip it up
		!cp /content/labelmap.txt /content/custom_model_lite
		!cp /content/labelmap.pbtxt /content/custom_model_lite
		!cp /content/models/mymodel/pipeline_file.config /content/custom_model_lite
		
		%cd /content
		!zip -r custom_model_lite.zip custom_model_lite
		
		from google.colab import files
		
		files.download('/content/custom_model_lite.zip')
	\end{lstlisting}
	
	The file \FILE{custom\_model\_lite.zip} containing the model will download into your Downloads folder. It's ready to be deployed on your device.
	
	%\subsubsection{Deploy model}
	TensorFlow Lite models can run on a wide variety of hardware, including PCs, embedded systems, and phones. So, the same object detection model can be run on Raspberry Pi or Arduino devices like Portenta H7.
	
